An observer in Hanoi (latitude of exactly $\phi = 21^\circ\,\mathrm{N}$) uses
an altazimuth telescope to observe Zeta Tauri (declination
$\delta = +21^\circ 08' 33''$) with a maximum azimuthal tracking speed of
$1.0^\circ/\mathrm{s}$. As an object transits close to the zenith, the
required tracking angular velocity increases sharply. If the required angular
velocity exceeds the motor's maximum speed, the telescope loses its target.
This problem is known as the ``Zenith blind spot'' or the ``Gimbal lock''. The
position of the star is defined by its azimuth ($A$), zenith distance ($z$),
and hour angle ($H$). Note that $A$ is measured Eastward from North and $H$ is
measured Westward from South.

\begin{description}
\item[(a)] (4 points) Show that for the telescope to be able to track the
  star, the required azimuthal angular velocity at meridian transit is given
  by $\omega_A = -\omega_\oplus \dfrac{\cos\delta}{\sin z}$ where
  $\omega_\oplus = 15^\circ/\mathrm{hr}$ is the Earth's sidereal rotation
  rate.
\item[(b)] (4 points) Calculate the minimum zenith distance
  ($z_{\mathrm{min}}$) reached by Zeta Tauri during its transit and the radius
  of the blind spot ($z_{\mathrm{limit}}$) in arcseconds. Show that Zeta Tauri
  falls into the blind spot of the telescope during its transit over Hanoi.
\item[(c)] (3 points) Assuming the star's path across the sky near the zenith
  can be approximated as a straight line, calculate the total time $\Delta t$
  (in seconds) that Zeta Tauri spends inside the telescope's blind spot.
\item[(d)] (4 points) Suppose a star passes exactly through the zenith; find
  the observer's latitudes $\phi_{\mathrm{max}}$ and $\phi_{\mathrm{min}}$ so
  that the sizes of the blind spots are maximal and minimal, respectively.
\end{description}
